% 第15批：《大学物理学学习指导》第5版自测题3填空3.2.1—3.2.8
% 印刷56／PDF物理64页。
\subsection{自测题3：光速、时间延缓和相对论密度填空}

\sourcetest{3.2.1}{56}{64}
\problemtext 恒星相对地球以速率$v$匀速运动；
恒星辐射的光子相对地球的速率大小是多少？
\workings
狭义相对论的真空光速不变，
\[
\boxed{u_\gamma=c}.
\]
\crosscheck 不能用伽利略相加写成$c+v$或$c-v$。

\sourcetest{3.2.2}{56}{64}
\problemtext $S'$沿$S$系$x$轴\emph{负方向}
以$0.5c$运动；$S'$原点O'沿$x'$正向发出光波。
$S$测量该光速大小。
\workings
相对运动速度$v=-0.5c$，
速度合成
\[
u=\frac{c-0.5c}{1-(0.5)}=c,
\]
故$\boxed{u=c}$。
\crosscheck 对光速使用洛伦兹速度合成；
不能得到经典的$0.5c$。

\sourcetest{3.2.3}{56}{64}
\problemtext 不稳定$\pi^+$介子的自身平均寿命为
$2.6\times10^{-8}$s，以$0.8c$相对实验室运动，
求实验室平均寿命。
\workings
运动钟时间延缓
\[
\boxed{\tau_{\rm lab}
=\frac{\tau_0}{\sqrt{1-0.8^2}}
=\frac{2.6\times10^{-8}}{0.6}
\simeq4.33\times10^{-8}\,\mathrm s}.
\]
\crosscheck 这是\emph{平均寿命}的参考系变换，
不是每个粒子实际存活时间都等于该值。

\sourcetest{3.2.4}{56}{64}
\problemtext 两位观察者相向接近的相对速度为$0.6c$，
其中O测得初始距离20m。
求O'自己的时钟从该初始时刻到相遇所历时间。
\workings
在O静止系，两者相遇需
$t_O=20/(0.6c)$。
O'时钟测量自身世界线上两个事件的固有时间，
\[
\boxed{\Delta\tau_{O'}=
t_O\sqrt{1-0.6^2}
=\frac{20(0.8)}{0.6c}
\simeq8.89\times10^{-8}\,\mathrm s}.
\]
\crosscheck 20m距离是由O系同时量得的初始距离；
不能直接套O'系长度20m而不考虑收缩。

\sourcetest{3.2.5}{56}{64}
\problemtext (1) 粒子的相对论动量大小何时等于相同速率下的
经典动量大小的2倍？
(2) 相对论动能何时等于自身静能？
\workings
(1) $p_{\rm rel}=\gamma m_0v$，
$p_{\rm classical}=m_0v$，要求$\gamma=2$。
(2) $K=(\gamma-1)m_0c^2=m_0c^2$，
也要求$\gamma=2$，所以两问速度相同：
\[
\boxed{v=c\sqrt{1-\tfrac1{2^2}}
=\frac{\sqrt3}{2}c}.
\]
\crosscheck 这里按静质量$m_0$定义经典动量，
不是沿用运动质量$m_{\rm rel}=\gamma m_0$后又多乘一个$\gamma$。

\sourcetest{3.2.6}{56}{64}
\problemtext 电子由静止加速到$0.6c$，求外力所做功，
以电子静质量$m_0$表达。
\workings
$\gamma=(1-0.6^2)^{-1/2}=1.25$，
\[
\boxed{W=\Delta K=(\gamma-1)m_0c^2
=\tfrac14m_0c^2}.
\]
\crosscheck 经典值$\tfrac12 m_0(0.6c)^2=0.18m_0c^2$，
与所求相对论值$0.25m_0c^2$不同。

\sourcetest{3.2.7}{56}{64}
\problemtext 观察者甲相对乙以$4c/5$运动，
甲携带静止长度$l$、截面积$S$、静止质量$m$的细棒，
细棒长度沿运动方向。
求甲、乙各自记录的棒的密度。
\workings
甲在棒自身静止系中测得
\[
\boxed{\rho_{\rm rest}=\frac{m}{Sl}}.
\]
乙测量的同时长度为
$l_B=l\sqrt{1-(4/5)^2}=3l/5$。
若使用\emph{现代静质量密度}的定义，
质量仍为$m$，故
\[
\boxed{\rho_{B,\rm restmass}=\frac{5m}{3Sl}}.
\]
但源书时代部分章节将“运动质量”
定义为$m_{\rm rel}=\gamma m=(5/3)m$；
若题中“密度”意为\emph{运动质量除同时体积}，
则乙得到
\[
\boxed{\rho_{B,\rm relativisticmass}
=\frac{(5/3)m}{(3/5)Sl}=\frac{25m}{9Sl}}.
\]
\crosscheck \textbf{源题术语约定必须独立二校：}
两种密度的数值不同，并非同一物理定义的两个解。
现代连续介质密度和相对论能动张量的定义又需考虑参考系切片；
本题只登记两种明确的教材简化约定，不虚称唯一无条件值。

\sourcetest{3.2.8}{56}{64}
\problemtext 速率为$0.99c$的电子静质量为
$9.11\times10^{-31}$kg。求相对论总能量（J）
及经典动能与相对论动能之比。
\workings
设$\gamma=1/\sqrt{1-0.99^2}\simeq7.0888$，
则
\[
\boxed{E=\gamma m_0c^2
\simeq(7.0888)(9.11\times10^{-31})(3.0\times10^8)^2
\simeq5.81\times10^{-13}\,\mathrm J}.
\]
经典动能与相对论动能分别是
$K_{\rm cl}=\tfrac12m_0v^2$、
$K_{\rm rel}=(\gamma-1)m_0c^2$，所以
\[
\boxed{\frac{K_{\rm cl}}{K_{\rm rel}}
=\frac{0.5(0.99)^2}{\gamma-1}
\simeq0.0805}.
\]
\crosscheck “总能量”包括静能，不是相对论动能本身。
